%% ----------------------------------------------------------------------------
%%  §1.10  Appendice: il principio di induzione
%% ----------------------------------------------------------------------------
\section{Appendice: il principio di induzione}
\label{sec:induzione}

Abbiamo utilizzato nel \cref{sec:potenze} la seguente affermazione sulla crescenza
della funzione potenza $x^{n}$:
\begin{equation}\label{eq:potenza-crescente-bis}
  0\le x_{1}<x_{2} \ \implica\ x_{1}^{n}<x_{2}^{n}, \qquad\perogni n\in\N,
\end{equation}
e ne abbiamo dato una dimostrazione per $n=2$ e $n=3$, concludendo con un «e così
via». Vogliamo ora rendere rigoroso quel «così via».

Supponiamo preliminarmente che la \eqref{eq:potenza-crescente-bis} valga per un certo
indice $n$ (ciò che dobbiamo provare è che valga per \emph{tutti} gli $n$; qui stiamo
supponendo che sia vera per \emph{qualche} $n$: ad esempio, per $n=1$ la
\eqref{eq:potenza-crescente-bis} è banalmente vera!). Perciò supponiamo che valgano le
disuguaglianze $0\le x_{1}<x_{2}$, $x_{1}^{n}<x_{2}^{n}$. Otteniamo, moltiplicando la
prima per $x_{1}^{n}\ge0$ e la seconda per $x_{2}>0$:
\begin{equation}\label{eq:passo-potenza}
  x_{1}^{n+1}=x_{1}\,x_{1}^{n}\le x_{2}\,x_{1}^{n}<x_{2}\,x_{2}^{n}=x_{2}^{n+1}.
\end{equation}
Cioè abbiamo provato che, \emph{se} vale la \eqref{eq:potenza-crescente-bis} per un
certo indice $n$, \emph{allora} essa vale anche per l'indice successivo $n+1$. Ma allora
la \eqref{eq:potenza-crescente-bis} vale sempre, perché: sappiamo che la proposizione
vale per $n=1$ (lo verifichiamo banalmente); per quanto sopra detto essa vale anche per
l'indice successivo, cioè $n=2$; ancora, sempre per lo stesso motivo, vale per il
successivo $n=3$, e così via\ldots{} Possiamo raggiungere con questo argomento
qualsiasi naturale $n$. Formuliamo in generale il seguente

\begin{teorema}[Principio di induzione]\label{teo:induzione}\index{principio di induzione}
Supponiamo che una proposizione $\mathcal{P}(n)$ dipendente da un indice $n\in\N$ sia
vera per $n=1$ e che inoltre, supposta vera per $n$, sia vera anche per il successivo
$n+1$. Allora la proposizione è vera per ogni $n\in\N$.
\end{teorema}

\begin{schema}
\Ipotesi{(a) \emph{Base:} $\mathcal{P}(1)$ è vera. \\
(b) \emph{Passo induttivo:} per ogni $n\in\N$, $\mathcal{P}(n)\implica\mathcal{P}(n+1)$.}
\Tesi{$\mathcal{P}(n)$ è vera per ogni $n\in\N$.}
\end{schema}

\begin{strategia}
Tutto sta nel dare un significato preciso alla frase «sommando ripetutamente $1$» con
cui abbiamo definito $\N$ nel \cref{sec:nzq}: $\N$ è il \emph{più piccolo}
sottoinsieme di $\R$ che contiene $1$ e che, se contiene $n$, contiene anche $n+1$.
Basta allora osservare che l'insieme dei naturali per cui $\mathcal{P}(n)$ è vera ha
queste due proprietà: essendo contenuto in $\N$ e non potendo essere più piccolo di
$\N$, coincide con $\N$.
\end{strategia}

\begin{dimostrazione}
Diciamo che un insieme $S\subseteq\R$ è \emph{induttivo} se $1\in S$ e se $n\in S$
implica $n+1\in S$. L'insieme $\R$ è induttivo, e l'intersezione di una famiglia
qualsiasi di insiemi induttivi è ancora un insieme induttivo (se $1$ e $n$ stanno in
tutti, vi stanno anche $1$ e $n+1$). La definizione di $\N$ come insieme dei numeri che
si ottengono da $1$ «sommando ripetutamente $1$» si rende rigorosa ponendo $\N$ uguale
all'intersezione di tutti i sottoinsiemi induttivi di $\R$: così $\N$ è induttivo ed è
contenuto in ogni insieme induttivo.

Sia ora $S=\{n\in\N : \mathcal{P}(n)\text{ è vera}\}$. Per l'ipotesi (a) è $1\in S$.
Se $n\in S$, allora $\mathcal{P}(n)$ è vera e, per l'ipotesi (b), lo è anche
$\mathcal{P}(n+1)$; inoltre $n+1\in\N$, perché $\N$ è induttivo: dunque $n+1\in S$.
Quindi $S$ è induttivo, e perciò $\N\subseteq S$; poiché per costruzione
$S\subseteq\N$, risulta $S=\N$, cioè $\mathcal{P}(n)$ è vera per ogni $n\in\N$.
\end{dimostrazione}

\begin{osservazione}[Il principio di induzione è un assioma o un teorema?]\label{oss:peano}
Nel nostro approccio, che parte dagli assiomi dei numeri reali, il principio di
induzione è un teorema, la cui dimostrazione si riduce a poche righe una volta precisato
che cosa significhi \emph{ripetutamente}. Nell'assiomatizzazione dei numeri naturali
proposta da Giuseppe Peano nel 1889, che non presuppone $\R$, si parte invece da un
insieme $\N$, da un suo elemento $1$ e da una funzione $\sigma$, il «successore», con
gli assiomi:
\begin{enumerate}[label=\textup{(P\arabic*)}, leftmargin=3em]
  \item $1\in\N$;
  \item se $n\in\N$, allora $\sigma(n)\in\N$;
  \item $1$ non è il successore di alcun numero: $\sigma(n)\neq1$ per ogni $n\in\N$;
  \item numeri con lo stesso successore sono uguali: se $\sigma(n)=\sigma(m)$, allora
  $n=m$;
  \item se $S\subseteq\N$ contiene $1$ e contiene $\sigma(n)$ ogni volta che contiene
  $n$, allora $S=\N$.
\end{enumerate}
Per Peano l'induzione è dunque un assioma, il quinto e ultimo. Il nostro $\N\subseteq\R$,
con $\sigma(n)=n+1$, li soddisfa tutti. (P1) e (P2) dicono che $\N$ è induttivo; (P4) è
la legge di semplificazione (\cref{prop:semplificazione-somma}); (P5) è il cuore della
dimostrazione del \cref{teo:induzione}, perché un tale $S$ è induttivo e quindi contiene
$\N$. Per (P3) basta osservare che l'insieme $\{x\in\R : x>0\}$ è induttivo (se $x>0$,
allora $x+1>1>0$), e quindi contiene $\N$: ogni naturale è positivo, mentre $n+1=1$
darebbe $n=0$. I naturali «dentro» $\R$ sono, a tutti gli effetti, i naturali di Peano.
\end{osservazione}

\begin{notastorica}[Peano e l'aritmetica «nova methodo»]
Giuseppe Peano (1858--1932), professore all'Università di Torino, pubblicò nel 1889 un
opuscolo in latino, \emph{Arithmetices principia, nova methodo exposita}, in cui
l'aritmetica è dedotta da pochi assiomi scritti per intero in simboli logici. Gli
assiomi erano nove, ma quattro riguardavano soltanto le proprietà dell'uguaglianza; i
cinque rimanenti sono quelli dell'\cref{oss:peano}, e il primo numero era, come in questo
libro, l'$1$: solo più tardi, nel \emph{Formulario mathematico}, Peano fece cominciare i
naturali dallo zero. Un anno prima Richard Dedekind, in \emph{Was sind und was sollen die
Zahlen?} (1888), aveva caratterizzato i naturali in modo equivalente, con il linguaggio
degli insiemi e delle funzioni, e Peano ne cita il lavoro nella prefazione. Somma e
prodotto si definiscono a partire dal successore per ricorrenza, come aveva già fatto
Hermann Grassmann nel 1861, e le loro proprietà si dimostrano per induzione. Dalla
simbologia di Peano viene anche un segno che usiamo dal \cref{sec:insiemi}: la lettera
$\varepsilon$, iniziale della parola greca \emph{esti} («è»), con cui Peano scriveva
$a\,\varepsilon\,\N$ per dire «$a$ è un numero»; da essa è nato il simbolo $\in$.
\end{notastorica}

\begin{idea}[La fila di tessere del domino]
L'immagine classica: una fila infinita di tessere del domino in piedi. La base (a) dice
che facciamo cadere la prima tessera; il passo induttivo (b) dice che ogni tessera, se
cade, urta la successiva e la fa cadere. Conclusione: cadono tutte. Si noti che
\emph{entrambe} le condizioni sono necessarie: senza (a) nessuna tessera cade mai;
senza (b) la prima cade e la seconda può restare in piedi. Si noti anche che nel passo
induttivo \emph{non} si afferma che $\mathcal{P}(n)$ è vera: si afferma solo che
\emph{se} lo fosse, lo sarebbe anche $\mathcal{P}(n+1)$. È un'implicazione, non
un'asserzione.
\end{idea}

\begin{proposizione}\label{prop:induzione-potenza}
La \eqref{eq:potenza-crescente-bis} vale per ogni $n\in\N$.
\end{proposizione}

\begin{schema}
\Ipotesi{$n\in\N$; \ $0\le x_{1}<x_{2}$.}
\Tesi{$x_{1}^{n}<x_{2}^{n}$.}
\end{schema}

\begin{strategia}
Si applica il \cref{teo:induzione} alla proposizione $\mathcal{P}(n)$ «per ogni
$0\le x_{1}<x_{2}$ risulta $x_{1}^{n}<x_{2}^{n}$». La base è l'ipotesi stessa; il passo
induttivo è il calcolo \eqref{eq:passo-potenza} già svolto: si tratta solo di
riconoscerlo come tale.
\end{strategia}

\begin{dimostrazione}
Sia $\mathcal{P}(n)$ la proposizione «per ogni $0\le x_{1}<x_{2}$ risulta
$x_{1}^{n}<x_{2}^{n}$». $\mathcal{P}(1)$ è vera (afferma che $x_{1}<x_{2}$). Il passo
induttivo è esattamente la \eqref{eq:passo-potenza}. Per il principio di induzione,
$\mathcal{P}(n)$ è vera per ogni $n$.
\end{dimostrazione}

Per chiarire meglio il principio consideriamo altri esempi: alcune formule classiche
che useremo più avanti, in particolare nella teoria delle serie e nella definizione
dell'integrale.

\begin{proposizione}[Somma dei primi $n$ numeri naturali]\label{prop:gauss}
Per ogni $n\in\N$,
\begin{equation}\label{eq:gauss}
  1+2+3+\dots+(n-1)+n=\frac{n(n+1)}{2}.
\end{equation}
\end{proposizione}

\begin{schema}
\Ipotesi{$n\in\N$.}
\Tesi{$1+2+\dots+n=\dfrac{n(n+1)}{2}$.}
\end{schema}

\begin{strategia}
In una dimostrazione per induzione di una \emph{formula} il passo induttivo consiste
sempre nel partire dal membro sinistro con $n+1$, isolare il membro sinistro con $n$,
sostituirlo con il membro destro con $n$ (ipotesi induttiva) e fare i conti fino a
ottenere il membro destro con $n+1$. Qui: $1+\dots+n+(n+1)=\frac{n(n+1)}{2}+(n+1)$, e
basta raccogliere $(n+1)$.
\end{strategia}

\begin{dimostrazione}
La formula è vera per $n=1$; infatti si ha l'identità $1=(1\cdot2)/2$. Supponiamo vera
la \eqref{eq:gauss} per un certo $n$ e dimostriamo la formula analoga con l'indice
$n+1$ al posto di $n$. Per ottenere ciò, è naturale sommare a entrambi i membri il
numero $n+1$:
\begin{equation}\label{eq:gauss-passo}
\begin{split}
  1+2+3+\dots+n&+(n+1)=\frac{n(n+1)}{2}+(n+1)\\
  &=\frac{n(n+1)+2(n+1)}{2}=\frac{(n+1)(n+2)}{2}.
\end{split}
\end{equation}
Abbiamo ottenuto ciò che volevamo: l'ultimo membro è esattamente il secondo membro della
\eqref{eq:gauss} con $n+1$ al posto di $n$. Quindi la \eqref{eq:gauss} risulta vera per
ogni $n\in\N$.
\end{dimostrazione}

\begin{notastorica}[Il piccolo Gauss]
Alla formula \eqref{eq:gauss} è legato un celebre aneddoto su Carl Friedrich Gauss
(1777--1855), raccontato dal suo primo biografo, Wolfgang Sartorius von Waltershausen
(1856): il maestro Büttner avrebbe assegnato alla classe la somma di una lunga
progressione aritmetica, e Gauss, che aveva allora circa nove anni, avrebbe posato quasi
subito sul tavolo la sua lavagnetta con il risultato esatto. Nella versione divenuta
popolare, fissata solo molto più tardi, i numeri da sommare sono quelli da $1$ a $100$ e
il risultato è $5050$. L'idea: si
accoppiano il primo e l'ultimo numero ($1+100=101$), il secondo e il penultimo
($2+99=101$), e così via: $50$ coppie da $101$. La dimostrazione per induzione che
abbiamo dato è meno brillante ma ha un pregio: funziona sempre, anche quando non si ha
un'idea geniale --- purché si conosca già la formula da provare. Il principio di
induzione \emph{verifica}, non \emph{scopre}.
\end{notastorica}

\begin{marginfigure}
\centering
\begin{tikzpicture}[scale=0.3]
  \foreach \i in {1,...,6} \foreach \j in {1,...,\i} \fill[blu!70] (\j-1,\i-1) rectangle ++(0.9,0.9);
  \foreach \i in {1,...,6} \foreach \j in {\i,...,6} \fill[rame!50] (\j,\i-1) rectangle ++(0.9,0.9);
  \draw[grigio, thin] (0,0) rectangle (7,6);
  \node[font=\scriptsize, below] at (3.5,-0.1) {$n+1$};
  \node[font=\scriptsize, left] at (-0.1,3) {$n$};
\end{tikzpicture}
\caption{Una «dimostrazione senza parole» della \eqref{eq:gauss}: due copie di
$1+2+\dots+n$ riempiono un rettangolo $n\times(n+1)$.}
\end{marginfigure}

Un'altra applicazione del principio di induzione è la seguente disuguaglianza, che
sarà fondamentale nello studio delle successioni.

\begin{teorema}[Disuguaglianza di Bernoulli]\label{teo:bernoulli}\index{disuguaglianza di Bernoulli}
Per ogni numero reale $x\ge-1$ e per ogni naturale $n$ risulta
\begin{equation}\label{eq:bernoulli}
  (1+x)^{n}\ge1+nx .
\end{equation}
\end{teorema}

\begin{schema}
\Ipotesi{$x\in\R$, $x\ge-1$; $n\in\N$.}
\Tesi{$(1+x)^{n}\ge1+nx$.}
\end{schema}

\begin{strategia}
Induzione su $n$ con $x$ fissato. Il passo induttivo richiede di moltiplicare
l'ipotesi $(1+x)^{n}\ge1+nx$ per $1+x$: ed è qui che serve l'ipotesi $x\ge-1$, perché
$1+x\ge0$ e la disuguaglianza non cambia verso (\cref{prop:ordine-prodotto}). Poi si
sviluppa il prodotto e si butta via il termine $nx^{2}\ge0$: è il tipico passaggio in
cui si \emph{indebolisce} una stima per ottenere la forma voluta.
\end{strategia}

\begin{dimostrazione}
Per $n=1$ la proposizione è vera (con il segno $=$). Supponiamo vera la
\eqref{eq:bernoulli} per un numero naturale $n$ e moltiplichiamo entrambi i membri per
$1+x$, che è una quantità maggiore o uguale a zero:
\begin{equation}\label{eq:bernoulli-passo}
  (1+x)^{n+1}\ge(1+nx)(1+x)=1+x+nx+nx^{2}\ge1+(n+1)x,
\end{equation}
dove nell'ultimo passaggio abbiamo trascurato il termine $nx^{2}\ge0$. Abbiamo ottenuto
la proposizione con $n+1$ al posto di $n$. Perciò, in base al principio di induzione, la
\eqref{eq:bernoulli} è provata.
\end{dimostrazione}

\begin{marginfigure}
\centering
\begin{tikzpicture}
\begin{axis}[width=5.2cm, height=4.4cm, xmin=-1.3, xmax=1.3, ymin=-2.5, ymax=5,
  xtick={-1,1}, ytick={1}, samples=100]
  \addplot[curva, domain=-1.25:1.25] {(1+x)^4};
  \addplot[curva2, domain=-1.25:1.25] {1+4*x};
  \node[blu, font=\scriptsize, anchor=east] at (axis cs:-0.12,3.9) {$(1+x)^{4}$};
  \node[rame, font=\scriptsize, anchor=west] at (axis cs:0.3,0.8) {$1+4x$};
  \addplot[only marks, mark=*, mark size=1.2pt, grigio] coordinates {(0,1)};
\end{axis}
\end{tikzpicture}
\caption{La disuguaglianza di Bernoulli per $n=4$: la curva $(1+x)^{4}$ sta sopra la
retta $1+4x$ (sua tangente in $x=0$) per ogni $x\ge-1$.}
\end{marginfigure}

\begin{osservazione}
Che cosa succede se $x<-1$? Per $n=2$ la \eqref{eq:bernoulli} vale ancora per ogni
$x$ (equivale a $x^{2}\ge0$), ma per $n=3$ e $x=-4$ si ha
$(1-4)^{3}=-27<-11=1+3\cdot(-4)$: la disuguaglianza è falsa. L'ipotesi $x\ge-1$ non è
dunque un eccesso di prudenza. Il lettore controlli, ripercorrendo la dimostrazione,
qual è esattamente il passaggio che fallisce quando $1+x<0$.
\end{osservazione}

Utilizzando il principio di induzione dimostriamo la formula che esprime la
\emph{somma di una progressione geometrica} di ragione $x\neq1$.

\begin{proposizione}[Somma della progressione geometrica]\label{prop:geometrica}\index{progressione geometrica}
Per ogni $x\in\R$, $x\neq1$, e per ogni $n\in\N$,
\begin{equation}\label{eq:geometrica}
  1+x+x^{2}+\dots+x^{n}=\frac{1-x^{n+1}}{1-x}.
\end{equation}
\end{proposizione}

\begin{schema}
\Ipotesi{$x\in\R$, \ $x\neq1$; \ $n\in\N$.}
\Tesi{$1+x+\dots+x^{n}=\dfrac{1-x^{n+1}}{1-x}$.}
\end{schema}

\begin{strategia}
Induzione su $n$, con $x$ fissato, secondo lo schema della \cref{prop:gauss}: si
aggiunge $x^{n+1}$ ai due membri dell'ipotesi induttiva e si riduce allo stesso
denominatore $1-x$, che non è nullo perché $x\neq1$. I termini $\pm x^{n+1}$ si
cancellano e resta il secondo membro con $n+1$ al posto di $n$.
\end{strategia}

\begin{dimostrazione}
Per $n=1$, a secondo membro abbiamo
\[
  \frac{1-x^{2}}{1-x}=\frac{(1-x)(1+x)}{1-x}=1+x;
\]
quindi la \eqref{eq:geometrica} è vera per $n=1$. Supponendo verificata la
\eqref{eq:geometrica} per un certo $n$, sommiamo a entrambi i membri il termine
$x^{n+1}$:
\begin{equation}\label{eq:geometrica-passo}
\begin{split}
  1+x+x^{2}+\dots+x^{n}+x^{n+1}&=\frac{1-x^{n+1}}{1-x}+x^{n+1}\\
  &=\frac{1-x^{n+1}+x^{n+1}-x^{n+2}}{1-x}=\frac{1-x^{n+2}}{1-x}.
\end{split}
\end{equation}
Abbiamo ottenuto la proposizione con $n+1$ al posto di $n$. Quindi, in base al
principio di induzione, la formula \eqref{eq:geometrica} è dimostrata.
\end{dimostrazione}

\begin{osservazione}[Una dimostrazione diretta]
La \eqref{eq:geometrica} si può anche ottenere senza induzione: posto
$S=1+x+\dots+x^{n}$, si ha $xS=x+x^{2}+\dots+x^{n+1}$, e sottraendo
$S-xS=1-x^{n+1}$ (tutti i termini intermedi si cancellano), da cui $S(1-x)=1-x^{n+1}$.
Ma anche questa dimostrazione nasconde un'induzione: il «si cancellano» per una somma
di $n$ termini è esso stesso un'affermazione che, a rigore, va provata per induzione.
\end{osservazione}

Proviamo mediante il principio di induzione una formula, analoga alla \eqref{eq:gauss},
che verrà utilizzata introducendo gli integrali definiti.

\needspace{16\baselineskip}
\begin{proposizione}[Somma dei primi $n$ quadrati]\label{prop:quadrati}
Per ogni $n\in\N$,
\begin{equation}\label{eq:quadrati}
  1^{2}+2^{2}+3^{2}+\dots+n^{2}=\frac{n(n+1)(2n+1)}{6}.
\end{equation}
\end{proposizione}

\begin{schema}
\Ipotesi{$n\in\N$.}
\Tesi{$1^{2}+2^{2}+\dots+n^{2}=\dfrac{n(n+1)(2n+1)}{6}$.}
\end{schema}

\begin{strategia}
Lo schema è quello della \cref{prop:gauss}, ma i conti sono più lunghi e conviene
organizzarli: dopo aver aggiunto $(n+1)^{2}$ si raccoglie il fattore comune $(n+1)$, si
riduce a un unico denominatore, e si confronta il polinomio di secondo grado ottenuto
con quello che compare nel membro destro con $n+1$ al posto di $n$. Due polinomi sono
uguali se hanno gli stessi coefficienti: è un controllo meccanico.
\end{strategia}

\begin{dimostrazione}
La formula \eqref{eq:quadrati} è vera per $n=1$; risulta infatti $1^{2}=(1\cdot2\cdot3)/6$.
Supponiamo che la \eqref{eq:quadrati} sia verificata per un indice $n$ e deduciamo da
essa la formula analoga con $n+1$ al posto di $n$. Abbiamo
\begin{equation}\label{eq:quadrati-passo}
\begin{split}
  1^{2}+2^{2}&+\dots+n^{2}+(n+1)^{2}
  =\bigl[1^{2}+2^{2}+\dots+n^{2}\bigr]+(n+1)^{2}\\
  &=\frac{n(n+1)(2n+1)}{6}+(n+1)^{2}\\
  &=(n+1)\Bigl[\frac{n(2n+1)}{6}+(n+1)\Bigr]\\
  &=(n+1)\,\frac{2n^{2}+n+6n+6}{6}
   =(n+1)\,\frac{2n^{2}+7n+6}{6}.
\end{split}
\end{equation}
Osservando che la \eqref{eq:quadrati} con $n+1$ al posto di $n$ ha a secondo membro
l'espressione
\begin{equation}\label{eq:quadrati-target}
  \frac{(n+1)(n+2)\bigl[2(n+1)+1\bigr]}{6}=\frac{(n+1)(n+2)(2n+3)}{6},
\end{equation}
rimane soltanto da osservare che
\[
  (n+1)\,\frac{2n^{2}+7n+6}{6}=\frac{(n+1)(n+2)(2n+3)}{6},
\]
che è una relazione soddisfatta perché, da verifica diretta, risulta
$2n^{2}+7n+6=(n+2)(2n+3)$.
\end{dimostrazione}

\begin{proposizione}[Una somma telescopica]\label{prop:telescopica}
Per ogni $n\in\N$,
\begin{equation}\label{eq:telescopica}
  \frac{1}{1\cdot2}+\frac{1}{2\cdot3}+\dots+\frac{1}{n(n+1)}=\frac{n}{n+1}.
\end{equation}
\end{proposizione}

\begin{schema}
\Ipotesi{$n\in\N$.}
\Tesi{$\dfrac{1}{1\cdot2}+\dots+\dfrac{1}{n(n+1)}=\dfrac{n}{n+1}$.}
\end{schema}

\begin{strategia}
Ancora lo schema consueto: si aggiunge il termine $\frac{1}{(n+1)(n+2)}$ ai due membri
dell'ipotesi induttiva, si riduce allo stesso denominatore e si riconosce al numeratore
il quadrato $(n+1)^{2}$, che si semplifica con il fattore $n+1$ del denominatore. Una
dimostrazione diretta, senza induzione, è indicata nell'osservazione che segue.
\end{strategia}

\begin{dimostrazione}
Da verifica diretta la formula \eqref{eq:telescopica} è soddisfatta per $n=1$.
Supponendo poi che la \eqref{eq:telescopica} sia verificata per un indice $n$ generico,
otteniamo
\begin{equation}\label{eq:telescopica-passo}
\begin{split}
  \frac{1}{1\cdot2}&+\dots+\frac{1}{n(n+1)}+\frac{1}{(n+1)(n+2)}
  =\frac{n}{n+1}+\frac{1}{(n+1)(n+2)}\\
  &=\frac{n(n+2)+1}{(n+1)(n+2)}=\frac{n^{2}+2n+1}{(n+1)(n+2)}\\
  &=\frac{(n+1)^{2}}{(n+1)(n+2)}=\frac{n+1}{n+2},
\end{split}
\end{equation}
che corrisponde appunto alla \eqref{eq:telescopica} con $n+1$ al posto di $n$.
\end{dimostrazione}

\begin{osservazione}
Il nome \emph{telescopica} viene dall'identità $\frac{1}{k(k+1)}=\frac1k-\frac{1}{k+1}$:
sommando per $k=1,\dots,n$ i termini intermedi si elidono a coppie, come i tubi di un
cannocchiale che rientrano l'uno nell'altro, e resta $1-\frac{1}{n+1}=\frac{n}{n+1}$.
Ritroveremo questa idea nello studio delle serie.
\end{osservazione}

\begin{python}[Verificare le formule con \texttt{sympy} e per forza bruta]
import sympy as sp
n, k, x = sp.symbols("n k x", positive=True, integer=True)

# sympy calcola le somme in forma chiusa: confrontiamo con le nostre
print(sp.factor(sp.summation(k, (k, 1, n))))          # Gauss
print(sp.factor(sp.summation(k**2, (k, 1, n))))       # quadrati
s = sp.summation(1/(k*(k+1)), (k, 1, n))
print(sp.simplify(s))                                 # telescopica

# Il passo induttivo della somma dei quadrati: due polinomi uguali?
print(sp.expand(2*n**2 + 7*n + 6 - (n + 2)*(2*n + 3)) == 0)

# Forza bruta: controlliamo la formula di Gauss per n = 1, ..., 1000
ok = (sum(range(1, m + 1)) == m*(m + 1)//2 for m in range(1, 1001))
print(all(ok))
\end{python}
\begin{risultato}
n*(n + 1)/2
n*(n + 1)*(2*n + 1)/6
n/(n + 1)
True
True
\end{risultato}

\begin{attenzione}[Mille casi non sono una dimostrazione]
L'ultima riga del programma controlla la \eqref{eq:gauss} per $n$ fino a $1000$:
rassicurante, ma \emph{non} è una dimostrazione. L'affermazione «$n^{2}+n+41$ è primo
per ogni $n$» è vera per $n=0,1,\dots,39$ e falsa per $n=40$ (Eulero); l'affermazione
«$991n^{2}+1$ non è mai un quadrato perfetto» è vera per tutti gli $n$ fino a circa
$10^{28}$, e poi falsa. Solo il principio di induzione (o un'altra dimostrazione)
garantisce \emph{tutti} gli $n$.
\end{attenzione}

\begin{esercizio}
Dimostrare per induzione che un insieme con $n$ elementi ha esattamente $2^{n}$
sottoinsiemi.
\begin{soluzione}
Per $n=1$: $\{a\}$ ha i due sottoinsiemi $\emptyset$ e $\{a\}$. Supponiamo che ogni
insieme con $n$ elementi abbia $2^{n}$ sottoinsiemi e sia $S$ un insieme con $n+1$
elementi; fissiamo $s\in S$. I sottoinsiemi di $S$ si dividono in quelli che non
contengono $s$ (sono i sottoinsiemi di $S\setminus\{s\}$, che ha $n$ elementi: $2^{n}$)
e quelli che contengono $s$ (si ottengono dai precedenti aggiungendo $s$: ancora
$2^{n}$). In tutto $2^{n}+2^{n}=2^{n+1}$. \emph{Scelta logica:} per il passo induttivo
bisogna ricondurre un insieme di $n+1$ elementi a uno di $n$: lo si fa isolando un
elemento $s$ e dividendo i sottoinsiemi in due classi disgiunte, a seconda che
contengano $s$ oppure no; ciascuna classe è in corrispondenza biunivoca con i
sottoinsiemi di $S\setminus\{s\}$.
\end{soluzione}
\end{esercizio}

\begin{esercizio}
Dimostrare per induzione che $2^{n}>n$ per ogni $n\in\N$, e che $n!\ge2^{n-1}$, dove
$n!=1\cdot2\cdots n$ è il \emph{fattoriale} di $n$.
\begin{soluzione}
(a) Per $n=1$: $2^{1}=2>1$. Se $2^{n}>n$, allora $2^{n+1}=2\cdot2^{n}>2n=n+n\ge n+1$,
perché $n\ge1$. (b) Per $n=1$: $1!=1=2^{0}$. Se $n!\ge2^{n-1}$, allora
$(n+1)!=(n+1)\,n!\ge(n+1)\,2^{n-1}\ge2\cdot2^{n-1}=2^{n}$, perché $n+1\ge2$ e
$2^{n-1}>0$. \emph{Scelta logica:} in entrambi i casi il passo induttivo confronta il
fattore per cui si moltiplica il primo membro passando da $n$ a $n+1$ ($2$ per $2^{n}$,
$n+1$ per $n!$) con quello del secondo membro; come nella disuguaglianza di Bernoulli,
si conclude \emph{indebolendo} la stima ($n\ge1$, $n+1\ge2$).
\end{soluzione}
\end{esercizio}
